\chapter{Robust envelopes under measurement slack}
\label{chap:robust-slack}

The robust-energy manuscript (\texttt{04\_robust\_energy}) is a
retirement candidate under the pivot, and this chapter is kept in a
single file so that retiring it is a deletion rather than a dissection.
Its subject is measurement slack in the embedding: the exposure bracket is
inflated by per-asset embedding radii, and the resulting interval is maximized
over a covariance box.
The pivot absorbs its motivation directly.
Slack is no longer an afterthought bolted onto an exact energy identity; it is
the \(\tau_i\) of \Cref{ass:random-common-map} and enters the transfer theorems
from the start.
What genuinely does not duplicate is \Cref{thm:box-corner}, a weights-only
maximization fact, and \Cref{thm:ball-inclusion}, which is the reconciliation
between the two geometries and should be promoted into
\Cref{chap:random-exposure} if the paper retires.

\section{Radii propagate through the embedding}

\begin{theorem}[Inner-product perturbation]
  \label{thm:inner-perturbation}
  \uses{ass:random-common-map}
  For mean embeddings perturbed by at most \(\varepsilon_i\) and
  \(\varepsilon_j\), the inner product moves by at most
  \(\varepsilon_i\lVert\mu_j\rVert+\varepsilon_j\lVert\mu_i\rVert
  +\varepsilon_i\varepsilon_j\).
\end{theorem}

\begin{proof}
  \uses{lem:polarization}
  Expand the perturbed inner product and bound each cross term.
\end{proof}

\begin{theorem}[Radii add on the metric]
  \label{thm:dist-bracket}
  \uses{ass:random-common-map}
  Under the same radii,
  \(\lvert\operatorname{dist}(\mu_i',\mu_j')-\operatorname{dist}(\mu_i,\mu_j)\rvert
  \le\varepsilon_i+\varepsilon_j\).
  Additivity on the metric, rather than on the squared quantity, is the same
  structural fact that makes the metric-level bracket of
  \Cref{cor:random-w2-metric} sharper than its squared counterpart.
\end{theorem}

\begin{proof}
  Two applications of the triangle inequality.
\end{proof}

\begin{theorem}[Asymmetric squared bracket]
  \label{thm:energy-bracket}
  \uses{thm:dist-bracket}
  Squaring the metric bracket gives an asymmetric interval, with the lower end
  truncated at zero.
  The asymmetry is not cosmetic: it is why the squared-cost formalization cannot
  simply inherit the metric statement.
\end{theorem}

\begin{proof}
  \uses{thm:dist-bracket}
  Square both ends of \Cref{thm:dist-bracket} and truncate.
\end{proof}

\begin{theorem}[Composed covariance bracket]
  \label{thm:covariance-bracket}
  \uses{ass:random-common-map, thm:energy-bracket}
  Under bi-Lipschitz constants stated in the true, unobserved distance and
  per-asset embedding radii, the covariance inner product is bracketed by
  polarization expressions evaluated at the radius-inflated and radius-deflated
  measured distance.
  The hypothesis that the bi-Lipschitz relation holds in the unobserved distance
  is the load-bearing one and is not testable from the measured distances.
\end{theorem}

\begin{proof}
  \uses{thm:energy-bracket, thm:inner-perturbation, lem:polarization}
  Compose \Cref{thm:energy-bracket} with the bi-Lipschitz relation and apply
  polarization at both ends.
\end{proof}

\section{Maximizing over the box}

\begin{theorem}[Long-only worst case is the upper corner]
  \label{thm:box-corner}
  \notready
  For nonnegative weights and a candidate covariance dominated entrywise by
  \(U\), the quadratic form is dominated by the form at \(U\); the supremum over
  the box is therefore attained at the upper corner.
  No positive-semidefiniteness of \(U\) is required, which is what makes the
  result usable on a bracket that is not itself a covariance matrix.
\end{theorem}

\begin{proof}
  Termwise domination against nonnegative weight products.
\end{proof}

\begin{theorem}[Robust portfolio envelope]
  \label{thm:robust-envelope}
  \uses{thm:covariance-bracket}
  For long-only weights the squared systematic risk lies between the quadratic
  forms of the entrywise lower and upper brackets built from the per-asset
  radii.
\end{theorem}

\begin{proof}
  \uses{thm:covariance-bracket, thm:box-corner}
  Apply \Cref{thm:box-corner} at both corners of the bracket from
  \Cref{thm:covariance-bracket}.
\end{proof}

\section{Reconciling the two geometries}

\begin{theorem}[Wasserstein balls sit inside energy balls]
  \label{thm:ball-inclusion}
  \uses{def:exposure-law, def:energy-negative-type}
  If \(W_2(\mu,\nu)\le\delta\) then \(\mathcal E^2(\mu,\nu)\le 2\delta\); the
  Wasserstein ball of radius \(\delta\) is therefore contained in the
  energy-distance ball of radius \(2\delta\).
  The inclusion is one-directional, and that direction is the pivot's
  justification: a \(W_2\) hypothesis is the stronger one, so results proved
  under it apply to the energy ball, while the converse fails.
  This node belongs with the trunk rather than with this chapter.
\end{theorem}

\begin{proof}
  \uses{def:energy-negative-type}
  Bound the energy distance by the transport cost through the negative-type
  representation.
\end{proof}
